% CODEX: wstaw bezpośrednio przed \section{Twierdzenie Hopfa--Rinowa}
% w 13-geodezyjne.tex. Pozostaw wcześniejsze przykłady geodezyjnych
% euklidesowych, sferycznych i hiperbolicznych.
\begin{przyklad}[Geodezyjne geometrii Sol i potencjał efektywny]
\label{prz:thur-dop-geodezyjne-sol}
Przyjmijmy metrykę $g_S$ z Przykładu~\ref{prz:thur-dop-metryki-nil-sol}
i niech $s\mapsto (x(s),y(s),t(s))$ będzie geodezyjną
w parametryzacji afinicznej. Jej lagranżjan kinetyczny jest równy
$L=\tfrac12(e^{-2t}\dot x^2+e^{2t}\dot y^2+\dot t^2)$.
Współrzędne $x,y$ nie występują w $L$, więc równania
Eulera--Lagrange'a dają dwie całki pierwsze
\begin{equation}\label{eq:thur-dop-sol-momenty}
 p_x=e^{-2t}\dot x,\qquad p_y=e^{2t}\dot y.
\end{equation}
Dzięki $\dot x=p_xe^{2t}$, $\dot y=p_ye^{-2t}$
stałość energii $E=L$ przyjmuje postać
\begin{equation}\label{eq:thur-dop-sol-energia}
 \tfrac12\dot t^2+V_{\mathrm{eff}}(t)=E,\qquad
 V_{\mathrm{eff}}(t)=\tfrac12(p_x^2e^{2t}+p_y^2e^{-2t}).
\end{equation}
Trzecie równanie geodezyjne jest równoważne
$\ddot t=-V_{\mathrm{eff}}'(t)
=-p_x^2e^{2t}+p_y^2e^{-2t}$.
W ten sposób trójwymiarowy układ redukuje się do ruchu
w jednowymiarowym potencjale i dwóch kwadratur dla $x,y$.
Jeżeli $p_xp_y\ne0$, potencjał dąży do nieskończoności
przy $t\to\pm\infty$ i ma jedno minimum przy
$e^{4t}=p_y^2/p_x^2$. Geodezyjna o skończonej energii
nie może wówczas uciec do nieskończoności w kierunku $t$.
Jest to twierdzenie o tych geodezyjnych, a nie zastępstwo
ogólnego twierdzenia o zupełności geometrii Sol.
\end{przyklad}

\begin{przyklad}[Rzuty geodezyjnych Nil na płaszczyznę]
\label{prz:thur-dop-geodezyjne-nil}
Dla $g_N$ lagranżjan wynosi
$L=\tfrac12(\dot x^2+\dot y^2+(\dot z-x\dot y)^2)$.
Współrzędna $z$ jest cykliczna, zatem
$w:=\dot z-x\dot y$ jest stałe. Równania
Eulera--Lagrange'a dla $x,y$ sprowadzają się do
\begin{equation}\label{eq:thur-dop-nil-geodezyjne}
 \ddot x=-w\dot y,\qquad \ddot y=w\dot x,
 \qquad \dot z=w+x\dot y.
\end{equation}
Rzeczywiście, równanie dla $y$ stwierdza stałość
$\dot y-xw$, a dla $x$ daje $\ddot x=-w\dot y$.
Wektor $(\dot x,\dot y)$ obraca się na płaszczyźnie
ze stałą prędkością kątową $w$ i ma stałą długość.
Jeśli $w=0$, rzut geodezyjnej jest prostą;
przy $w\ne0$ i niezerowej prędkości poziomej jest łukiem okręgu.
Ostatnie równanie wyznacza współrzędną pionową przez całkowanie.
Nie należy utożsamiać geodezyjnych Nil z krzywymi poziomymi:
warunek poziomości to $w=0$ i jest dodatkowym ograniczeniem.
\end{przyklad}
